\documentclass[conference]{IEEEtran}
\IEEEoverridecommandlockouts

\usepackage[utf8]{inputenc}
\usepackage[T1]{fontenc}
\usepackage[turkish]{babel}
\usepackage{graphicx}
\usepackage{amsmath,amssymb}
\usepackage{booktabs}
\usepackage{multirow}
\usepackage{float}
\usepackage{cite}
\usepackage{xcolor}
\usepackage[hidelinks]{hyperref}

\graphicspath{{./}}

% ─────────────────────────────────────────────────────────────────────────────
\title{EfficientNet-B0 Tabanlı Araba Gövde Tipi Sınıflandırması}

\author{
  \IEEEauthorblockN{Barışcan Özgen}
  \IEEEauthorblockA{
    Bilgisayar Mühendisliği Bölümü\\
    Kocaeli Üniversitesi\\
    Kocaeli, Türkiye\\
    230201057}
}

\begin{document}
\maketitle

% ─────────────────────────────────────────────────────────────────────────────
\begin{abstract}
Bu çalışmada, sekiz farklı araba gövde tipini (SUV, Van, Station Wagon, Micro,
Açık Tekerlekli, Sedan, Hatchback ve Pick-Up) sınıflandırmak amacıyla
EfficientNet-B0 mimarisi üzerine kurulu bir transfer öğrenme sistemi
geliştirilmiştir. Kaggle veri setleri ve DuckDuckGo görüntü arama motoru
aracılığıyla derlenen 6.400 görsellik özel bir veri seti oluşturulmuş;
CLIP (Contrastive Language-Image Pre-Training) modeli ile semantik kalite
filtresi uygulanmıştır. Veri seti \%70/\%15/\%15 oranında eğitim, doğrulama
ve test kümelerine ayrılmıştır. Adam optimizer, ReduceLROnPlateau zamanlayıcı
ve erken durdurma (Early Stopping) teknikleriyle eğitilen model, test kümesinde
\%93,54 doğruluk ve 0,9357 Macro F1 skoru elde etmiştir. Eğitilen model,
gerçek zamanlı tahmin yapabilen FastAPI tabanlı bir web arayüzüne entegre
edilmiştir.
\end{abstract}

\begin{IEEEkeywords}
Araba Gövde Tipi Sınıflandırma, EfficientNet-B0, Transfer Öğrenme,
CLIP, Derin Öğrenme, FastAPI, Görüntü Sınıflandırma
\end{IEEEkeywords}

% ─────────────────────────────────────────────────────────────────────────────
\section{Giriş}

Araç tipi tanıma, akıllı trafik yönetimi, otopark sistemleri, otonom sürüş
ve sigorta uygulamaları gibi birçok gerçek dünya probleminde kritik bir rol
oynamaktadır. Geleneksel görüntü işleme yöntemleri bu alanda sınırlı başarı
sergilemiş olsa da derin öğrenme tekniklerinin gelişmesiyle birlikte sınıflama
doğruluğu önemli ölçüde artmıştır \cite{lecun2015deep}.

Convolutional Neural Network (CNN) mimarileri, özellikle ImageNet gibi büyük
veri kümeleri üzerinde önceden eğitilmiş modeller, az sayıda etiketli veriyle
bile yüksek sınıflandırma doğruluğu elde edilmesine olanak sağlamaktadır
\cite{tan2019efficientnet}. Bu çalışmada, parametre verimliliği açısından
üstün bir performans sunan EfficientNet-B0 mimarisi seçilmiş; ImageNet
ağırlıkları kullanılarak transfer öğrenme uygulanmıştır.

Çalışmanın temel katkıları şöyle özetlenebilir:
\begin{itemize}
  \item Kaggle ve DuckDuckGo kaynaklarından derlenen, CLIP tabanlı semantik
        kalite filtresiyle temizlenen özgün bir veri seti oluşturulması.
  \item EfficientNet-B0 mimarisinin sekiz sınıflı araba gövde tipi
        sınıflandırmasına transfer öğrenme yöntemiyle uyarlanması.
  \item Eğitilmiş modelin hızlı çıkarım kapasiteli bir web arayüzüne
        entegrasyonu.
\end{itemize}

% ─────────────────────────────────────────────────────────────────────────────
\section{Veri Seti}

\subsection{Veri Toplama}

Veri toplama süreci birden fazla kaynaktan yürütülmüştür. İlk aşamada
Kaggle üzerinde yer alan \emph{Cars Body Type Cropped} \cite{boukhris2023cars}
ve \emph{Stanford Car Body Type Data} \cite{krause2013stanford} veri setleri
kullanılmıştır. Birinci veri seti SUV, Van, Sedan, Hatchback ve Pick-Up
sınıfları için temel görsel kaynağı oluştururken; ikinci veri seti
Station Wagon sınıfının tamamlanmasında kullanılmıştır.

Micro ve Station Wagon sınıfları için ek görsel ihtiyacı, \texttt{ddgs} Python
kütüphanesi aracılığıyla DuckDuckGo görüntü arama motoru üzerinden otomatik
indirme yöntemiyle karşılanmıştır. Station Wagon için ``station wagon car side
view'', ``estate car'', ``kombi araba'' gibi; Micro için ``smart fortwo'',
``fiat 500 side view'', ``microcar side view'' gibi arama terimleri
kullanılmıştır. F1 (Açık Tekerlekli) sınıfı için ise ``formula 1 car side
view'' ve ``open wheel racing car'' arama terimlerine ek olarak kişisel
koleksiyon görsellerinden faydalanılmıştır.

\subsection{Semantik Kalite Filtresi}

Ham görsel toplamada oluşabilecek gürültüyü (logo, çizim, iç mekan
fotoğrafları, alakasız nesneler vb.) gidermek amacıyla OpenAI'nin CLIP
\cite{radford2021learning} modeli (ViT-B/32 omurgası) kullanılmıştır. Her
görsel için ``a photo of a [sınıf] car'' metin sorgusuyla hesaplanan kosinüs
benzerlik skoru 0,25 eşiğinin altında kalan görseller veri setinden çıkarılmış
ve yerlerine yeni görseller eklenmiştir. Bu işlem, her sınıf tam 800 kaliteli
görsele ulaşana kadar tekrarlanmıştır.

\subsection{Veri Seti İstatistikleri}

Tablo \ref{tab:dataset}'de nihai veri setinin sınıf bazlı dağılımı
verilmiştir. Her sınıfta 800 görsel bulunmakta olup toplam 6.400 görsel
\%70/\%15/\%15 oranında eğitim/doğrulama/test kümelerine rastgele
bölünmüştür (sabit seed=42).

\begin{table}[h]
\caption{Veri Seti Dağılımı}
\label{tab:dataset}
\centering
\begin{tabular}{lcccc}
\toprule
\textbf{Sınıf} & \textbf{Toplam} & \textbf{Eğitim} & \textbf{Doğrulama} & \textbf{Test} \\
\midrule
SUV             & 800 & 560 & 120 & 120 \\
Van             & 800 & 560 & 120 & 120 \\
Station Wagon   & 800 & 560 & 120 & 120 \\
Micro           & 800 & 560 & 120 & 120 \\
Açık Tekerlekli & 800 & 560 & 120 & 120 \\
Sedan           & 800 & 560 & 120 & 120 \\
Hatchback       & 800 & 560 & 120 & 120 \\
Pick-Up         & 800 & 560 & 120 & 120 \\
\midrule
\textbf{Toplam} & \textbf{6400} & \textbf{4480} & \textbf{960} & \textbf{960} \\
\bottomrule
\end{tabular}
\end{table}

% ─────────────────────────────────────────────────────────────────────────────
\section{Yöntem}

\subsection{Ön İşleme ve Veri Çoğaltma}

Tüm görseller modelin beklediği 224$\times$224 piksel boyutuna yeniden
ölçeklendirilmiş ve ImageNet istatistikleriyle
($\mu$=[0.485, 0.456, 0.406], $\sigma$=[0.229, 0.224, 0.225])
normalize edilmiştir.

Eğitim kümesinde aşırı öğrenmeyi azaltmak ve modelin farklı koşullara
genelleme yeteneğini artırmak amacıyla veri çoğaltma (data augmentation)
uygulanmıştır. Uygulanan dönüşümler şunlardır:

\begin{itemize}
  \item \textbf{RandomHorizontalFlip} ($p=0,5$): Yatay çevirme
  \item \textbf{RandomRotation} ($\pm15°$): Rastgele döndürme
  \item \textbf{ColorJitter} (brightness=0,2; contrast=0,2;
        saturation=0,2; hue=0,05): Renk dönüşümleri
  \item \textbf{RandomPerspective} (distortion=0,2; $p=0,3$):
        Perspektif bozma
\end{itemize}

Doğrulama ve test kümelerine yalnızca yeniden ölçekleme ve normalizasyon
uygulanmıştır.

\subsection{Model Mimarisi}

EfficientNet-B0 \cite{tan2019efficientnet}, bileşik ölçekleme (compound
scaling) yöntemiyle derinlik, genişlik ve çözünürlük boyutlarını birlikte
optimize eden verimli bir CNN mimarisidir. Yaklaşık 5,3 milyon parametreye
ve ImageNet'te \%77,1 top-1 doğruluğa sahip olan bu model,
hesaplama maliyeti-doğruluk dengesi açısından üstün bir tercih sunmaktadır.

Bu çalışmada, ImageNet-1K ağırlıklarıyla önceden eğitilmiş EfficientNet-B0
modeli transfer öğrenme çerçevesinde kullanılmıştır. Orijinal 1.000 nöronlu
sınıflandırma katmanı, projemizin 8 sınıfına uygun
$\text{Linear}(1280 \rightarrow 8)$ katmanıyla değiştirilmiştir.
Modelin tüm parametreleri ince ayar (fine-tuning) sürecinde
eğitilebilir olarak bırakılmıştır. Nihai model boyutu yaklaşık 46 MB
olup proje sınırı olan 95 MB'nin çok altındadır.

\subsection{Eğitim Detayları}

Model, NVIDIA RTX 4070 Laptop GPU üzerinde PyTorch 2.x ile eğitilmiştir.
Kullanılan hiperparametreler Tablo \ref{tab:hyperparams}'de verilmiştir.

\begin{table}[h]
\caption{Hiperparametreler}
\label{tab:hyperparams}
\centering
\begin{tabular}{ll}
\toprule
\textbf{Hiperparametre} & \textbf{Değer} \\
\midrule
Görüntü boyutu      & $224 \times 224$ piksel \\
Batch boyutu        & 32 \\
Maksimum epoch      & 50 \\
Öğrenme oranı       & $10^{-4}$ \\
Ağırlık bozulması   & $10^{-4}$ \\
Kayıp fonksiyonu    & CrossEntropyLoss \\
Optimizer           & Adam \cite{kingma2014adam} \\
LR Zamanlayıcı      & ReduceLROnPlateau (factor=0,5; patience=3) \\
Erken Durdurma      & Val. kaybına göre (patience=7) \\
\bottomrule
\end{tabular}
\end{table}

Kayıp fonksiyonu olarak çok sınıflı sınıflandırma için standart olan
Cross-Entropy Loss kullanılmıştır. Öğrenme oranı, doğrulama kaybında
üst üste 3 epoch boyunca iyileşme gözlenmediğinde 0,5 çarpanıyla
azaltılmıştır. Doğrulama kaybında 7 epoch boyunca iyileşme
sağlanamadığında eğitim otomatik olarak durdurulmuştur.

% ─────────────────────────────────────────────────────────────────────────────
\section{Deneysel Sonuçlar}

\subsection{Eğitim Süreci}

Model, erken durdurma mekanizması devreye girdiğinde 18. epoch'ta eğitimi
tamamlamıştır. En iyi model ağırlıkları, 11. epoch'ta elde edilen
Val. Loss = 0,1345 ve Val. Acc = \%95,10 değerlerine karşılık gelen
kontrol noktasından alınmıştır.

Şekil \ref{fig:curves}'de eğitim/doğrulama
kayıp ve doğruluk eğrileri gösterilmiştir. Eğitim kaybının monoton biçimde
düşmesi ve doğrulama kaybının belirli bir değerde stabilize olması, modelin
veri setine uygun biçimde öğrendiğini göstermektedir. Eğitim ve doğrulama
doğrulukları arasındaki dar boşluk, ciddi bir aşırı öğrenme (overfitting)
olmadığının işaretidir.

\begin{figure*}[!t]
  \centering
  \includegraphics[width=0.48\textwidth]{training_loss.png}
  \hfill
  \includegraphics[width=0.48\textwidth]{training_accuracy.png}
  \caption{Sol: Eğitim ve Doğrulama Kayıp Eğrisi; Sağ: Eğitim ve Doğrulama Doğruluk Eğrisi (18 Epoch)}
  \label{fig:curves}
\end{figure*}

\subsection{Performans Metrikleri}

Model, eğitim sürecinde hiç kullanılmamış 960 görsellik test kümesi üzerinde
değerlendirilmiştir.
Tablo \ref{tab:results}'de sınıf bazlı Precision,
Recall ve F1 skorları verilmiştir.

\begin{table}[h]
\caption{Test Kümesi Sınıflandırma Sonuçları (support=120/sınıf)}
\label{tab:results}
\centering
\begin{tabular}{lccc}
\toprule
\textbf{Sınıf} & \textbf{Precision} & \textbf{Recall} & \textbf{F1} \\
\midrule
Açık Tekerlekli & 1,0000 & 1,0000 & 1,0000 \\
Hatchback       & 0,8254 & 0,8667 & 0,8455 \\
Micro           & 0,9237 & 0,9083 & 0,9160 \\
Pick-Up         & 0,9667 & 0,9667 & 0,9667 \\
Sedan           & 0,9720 & 0,8667 & 0,9163 \\
Station Wagon   & 0,8692 & 0,9417 & 0,9040 \\
SUV             & 0,9661 & 0,9500 & 0,9580 \\
Van             & 0,9752 & 0,9833 & 0,9793 \\
\midrule
\textbf{Macro Ort.}    & 0,9373 & 0,9354 & \textbf{0,9357} \\
\textbf{Weighted Ort.} & 0,9373 & 0,9354 & \textbf{0,9357} \\
\midrule
\multicolumn{3}{l}{\textbf{Test Doğruluğu}} & \textbf{\%93,54} \\
\bottomrule
\end{tabular}
\end{table}

\begin{figure*}[!t]
  \centering
  \includegraphics[width=0.33\textwidth]{confusion_matrix.png}
  \caption{Normalize Edilmiş Karışıklık Matrisi (Test Kümesi, $n$=960)}
  \label{fig:cm}
\end{figure*}

\subsection{Karışıklık Matrisi Analizi}

Şekil \ref{fig:cm}'de normalize edilmiş karışıklık matrisi gösterilmektedir.
Açık Tekerlekli (F1) sınıfı mükemmel bir \%100 sınıflandırma başarısı
sergilemiştir; bunun temel nedeni bu araç tipinin diğer sınıflardan belirgin
biçimde farklı olan açık kokpit ve tek koltuk tasarımıdır.

En düşük performans Hatchback sınıfında gözlemlenmiştir (F1=0,8455). Bu
durum, Hatchback ve Sedan araç tiplerinin benzer gövde hatları paylaşmasından
kaynaklanmaktadır; küçük bagaj açıklığı ve arka cam eğimi gibi ince
görsel farklılıklar modelin hata yapmasına yol açabilmektedir.
Station Wagon sınıfı da benzer nedenlerle Sedan ile zaman zaman karışmaktadır
(F1=0,9040).

% ─────────────────────────────────────────────────────────────────────────────
\section{Web Arayüzü}

Eğitilen model, FastAPI \cite{fastapi} framework'ü kullanılarak bir REST API
aracılığıyla sunulmuştur. Kullanıcı arayüzü saf HTML, CSS ve Vanilla
JavaScript ile geliştirilmiş olup herhangi bir frontend framework bağımlılığı
bulunmamaktadır.

Arayüz dört temel bileşenden oluşmaktadır:
\begin{enumerate}
  \item \textbf{Görüntü Yükleme Alanı:} Sürükle-bırak (drag-and-drop) ve
        dosya seçici (file picker) desteğiyle JPG, PNG ve WEBP
        formatlarını kabul etmektedir.
  \item \textbf{Görüntü Önizlemesi:} Yüklenen görsel anında
        ekranda gösterilmektedir.
  \item \textbf{Tahmin Butonu:} ``Tahmin Yap'' butonu ile POST /predict
        endpoint'ine görsel gönderilmektedir.
  \item \textbf{Sonuç Paneli:} Tahmin edilen sınıf etiketi, güven skoru ve
        tüm 8 sınıf için olasılık dağılımını gösteren yatay bir
        çubuk grafik (Chart.js) içermektedir.
\end{enumerate}

Modelin çıkarım süresi NVIDIA RTX 4070 Laptop GPU üzerinde görüntü başına
yaklaşık 50--100 ms olmaktadır; bu değer etkileşimli kullanım için
kabul edilebilir bir gecikme süresidir.

% ─────────────────────────────────────────────────────────────────────────────
\section{Tartışma ve Sonuç}

Bu çalışmada, EfficientNet-B0 tabanlı transfer öğrenme ile sekiz sınıflı
araba gövde tipi sınıflandırması başarıyla gerçekleştirilmiştir. Test
kümesinde elde edilen \%93,54 doğruluk ve 0,9357 Macro F1 skoru, modelin
görülmemiş verilere iyi bir genelleme yeteneğine sahip olduğunu
göstermektedir.

Veri toplama sürecinde CLIP tabanlı semantik filtre, veri kalitesini önemli
ölçüde artırmış; alakasız görsellerin eğitim ve test kümelerine karışmasının
önüne geçmiştir. Her sınıf için tam 800 kaliteli görsel sağlanması,
sınıflar arası dengeyi korumuştur.

Gelecek çalışmalar kapsamında şu iyileştirmeler değerlendirilebilir:
\begin{itemize}
  \item Hatchback sınıfına özgü ek veri toplama ile bu sınıftaki
        performansın artırılması.
  \item EfficientNet-B0'ın son birkaç bloğunun da ince ayar sürecine dahil
        edilerek sınıfa özgü özellik öğreniminin güçlendirilmesi.
  \item Test zamanı veri çoğaltma (Test Time Augmentation — TTA)
        yöntemiyle tahmin kararlılığının artırılması.
\end{itemize}

Sonuç olarak bu proje, uygun büyüklükte ve kalitede bir veri seti ile
transfer öğrenme yaklaşımının araba gövde tipi sınıflandırmasında
güvenilir ve hızlı sonuçlar üretebileceğini ortaya koymaktadır.

% ─────────────────────────────────────────────────────────────────────────────
\bibliographystyle{IEEEtran}
\bibliography{references}

\end{document}
